% from ega4-2-fr.tex
\oldpage[IV-2]{150}

En effet $f_\lambda$ est bijectif en vertu de (I, 3.5.2, 3.5.7 et 3.5.8);
d'autre part, il est fermé par (II, 6.1.10). On peut maintenant appliquer
(5.13.6) qui prouve (en vertu de l'hypothèse b)) que b) entraîne e) lorsque
$P(Z,z)$ est la propriété $(R_n)$ au point $z$. Le cas où $P(Z,z)$ est la
propriété d'être régulier (i.e. de posséder la propriété $(R_n)$ pour tout $n$)
en résulte trivialement. Enfin, compte tenu du critère de normalité de Serre
(5.8.6), b) entraîne encore e) lorsque $P(Z,z)$ est la propriété d'être normal
au point $z$, car cela équivaut à dire que $Z$ possède au point $z$ les
propriétés $(S_2)$ et $(R_1)$: on applique alors ce qui précède pour $n=1$,
et (6.7.2, (ii)) pour $n=2$.
\end{proof}

\begin{corollary}[6.7.8]
\label{IV.2.s6.7.8-fr}
Soit $k$ un corps; pour un $k$-préschéma localement noethérien $X$ et un
point $x\in X$, on désigne par $P(X,x)$ une des propriétés suivantes:
\begin{enumerate}
\item[(i)] $\operatorname{coprof}(\mathcal O_x)\leq n$.
\item[(ii)] $\mathcal O_x$ est un anneau de Cohen--Macaulay.
\item[(iii)] $X$ vérifie la propriété $(S_n)$ au point $x$.
\item[(iv)] $X$ est géométriquement régulier au point $x$.
\item[(v)] $X$ possède la propriété $(R_n)$ géométrique au point $x$.
\item[(vi)] $X$ est géométriquement normal au point $x$.
\item[(vii)] $X$ est géométriquement réduit (i.e. séparable) au point $x$.
\end{enumerate}
Soit $k'$ une extension de $k$; on suppose, soit que $k'$ est de type fini,
soit que $X$ est localement de type fini sur $k$, de sorte que
$X'=X\otimes_k k'$ est localement noethérien. Soit $x'$ un point de $X'$ dont
$x$ est la projection dans $X$. Alors, pour que $P(X,x)$ soit vraie, il faut
et il suffit que $P(X',x')$ le soit.
\end{corollary}

\begin{proof}
Cela a déjà été vu pour la propriété (vii) (4.6.11), et pour les propriétés
(i), (ii) et (iii) (6.7.1). Pour (iv), (v) et (vi), cela résulte de (6.7.7):
en effet, le fait que la condition est nécessaire résulte de l'équivalence des
critères a) et c), et de l'équivalence de c) et d) lorsque $X$ est localement
de type fini sur $k$. Pour voir que la condition est suffisante, soit $k''$ une
extension finie radicielle de $k$; on peut toujours considérer $k'$ et $k''$
comme sous-extensions d'une extension $K$ de $k$; posons
$X''=X\otimes_k k''$, $X_0=X\otimes_k K$, et notons que puisque $k''$ est une
extension radicielle de $k$, il n'y a qu'un seul point $x''$ de $X''$ au-dessus
de $x$ (I, 3.5.7 et 3.5.8). Soit alors $x_0$ un point quelconque de $X_0$
au-dessus de $x'$; si $P(X',x')$ est vraie, il en est de même de $P(X_0,x_0)$
en vertu de (6.7.7, c) et d)) (car si $k'$ est une extension de $k$ de type
fini, on peut supposer qu'il en est de même de $K$, et si $X$ est localement
de type fini sur $k$, $X'$ est localement de type fini sur $k'$). On déduit
alors de (6.7.4) que la propriété $Q(k'')$ est vraie (avec les notations de
(6.7.7)), donc $P(X,x)$ est vraie par (6.7.7, b)).
\end{proof}

\subsection{Morphismes réguliers, normaux, réduits, lisses}
\label{subsection:IV.6.8}

\begin{definition}[6.8.1]
\label{IV.2.s6.8.1-fr}
Soient $X$, $Y$ deux préschémas, $f:X\to Y$ un morphisme tel que la fibre
$f^{-1}(y)$ soit un préschéma localement noethérien pour tout $y\in Y$, $x$
un point de $X$. On dit respectivement que $f$ est un morphisme:
\begin{enumerate}
\item[(i)] de coprofondeur $\leq n$ au point $x$;
\item[(ii)] de Cohen--Macaulay au point $x$;
